These notes develop the Ricci flow arguments in
``The Entropy Formula for the Ricci Flow and its Geometric
Applications'' \cite{Perelman} and ``Ricci Flow with Surgery
on Three-Manifolds'' \cite{Perelman2}.  They include the construction of
Ricci flow with surgery and its long-time analysis.  In the simply-connected
case, finite extinction results of
\cite{Colding-Minicozzi,Colding-Minicozzi2,Perelman3}, together with surgery
existence, give the Poincar\'e Conjecture; the general long-time analysis gives
the Geometrization Conjecture.

\section{An overview of the Ricci flow approach to
$3$-manifold geometrization}
\label{overview0}

All manifolds and Riemannian metrics in this section are smooth.  We use the
notation of \cite{Perelman,Perelman2}: $R$ is scalar curvature, $\Ric$ is
Ricci curvature, and $|\Rm|$ is the largest absolute sectional curvature.
The inequality $\Rm\geq C$ means that all sectional curvatures at the
specified point or region are at least $C$.  We specialize to dimension three;
the geometrization conjecture is stated in Appendix \ref{geom}.

\subsection{The definition of Ricci flow, and some basic properties}

Let $M$ be a compact $3$-manifold and let $\{g(t)\}_{t\in [a,b]}$ be a
smooth family of Riemannian metrics.  It satisfies the Ricci flow equation if
\begin{equation}
\frac{\D g}{\D t}(t) \: = \:
-\:2\:\Ric(g(t))
\end{equation}
for every $t\in[a,b]$.  Hamilton showed in \cite{Hamiltonnnnn} that for
every metric $g_0$ there is a unique maximal solution on $[0,T)$ with
$g(0)=g_0$, where $T\in(0,\infty]$ and, if $T<\infty$, the curvature is
unbounded as $t\ra T$.  The time $T$ is the blow-up time.  For
$g_0=r_0^2g_{S^3}$, the shrinking round solution is
$g(t)=(r_0^2-4t)g_{S^3}$ and $T=r_0^2/4$.

If $g_0$ has positive Ricci curvature, Hamilton proved in \cite{Hamilton}
that the solution becomes round and shrinks to a point in finite time.  In
general, a simply-connected flow can develop a neckpinch before extinction.
A neckpinch is modeled on $(-c,c)\times S^2$, where $S^2$-fibers shrink to
points.  Surgery removes a neighborhood diffeomorphic to
$(-c',c')\times S^2$ and caps the boundary components with $3$-balls.  The
resulting smooth manifold is used as the initial manifold for the next Ricci
flow; repeated surgery gives Ricci flow with surgery.

\subsection{A rough outline of the Ricci flow proof of the Poincar\'e
Conjecture}

Let $M$ be a compact orientable $3$-manifold with arbitrary metric $g_0$, and
let $g(\cdot)$ be the Ricci flow on $[0,T)$.  If $T<\infty$, let
\[
\Omega=\{x\in M:\lim_{t\rightarrow T^-}R(x,t)\text{ exists and is finite}\}.
\]
The singular set is $M-\Omega$.

\begin{proposition}[existence for Riemannian metrics]
  \label{prop:kl-overview-ricci-flow-approach-manifold-geometrization-existence-riemannian-metrics}
  \source{kleiner-lott}
  \dcref{Claim 1}
  \group{Kleiner--Lott Overview}
 \cite{Perelman2} \label{introclaim1}
The set $\Omega$ is open and, as $t\ra T$, $g(\cdot)$ converges smoothly on
compact subsets of $\Omega$ to a Riemannian metric $\overline g$.  There is a
geometrically defined neighborhood $U$ of $M-\Omega$ whose connected
components are either
\[
\begin{array}{ll}
\mathrm{A}.&\text{compact and diffeomorphic to }S^1\times S^2,
S^1\times_{\Z_2}S^2\text{, or }S^3/\Gamma,\\
&\Gamma\text{ a finite subgroup of }\SO(4)\text{ acting freely and
isometrically on the round }S^3,
\end{array}
\]
or
\[
\begin{array}{ll}
\mathrm{B}.&\text{noncompact and diffeomorphic to }\R\times S^2,
\R^3\text{, or }\R\times_{\Z_2}S^2\text{ over }\R P^2.
\end{array}
\]
Here the generator of $\Z_2$ acts on $S^1$ by complex conjugation and on
$S^2$ by the antipodal map, so
$S^1\times_{\Z_2}S^2\cong\R P^3\#\R P^3$.  In Case B, each component
meets $\Omega$ in geometrically controlled collar regions diffeomorphic to
$\R\times S^2$.
\end{proposition}

\begin{proposition}[existence for Ricci flow with surgery]
  \label{prop:kl-overview-ricci-flow-approach-manifold-geometrization-existence-ricci-flow-surgery}
  \source{kleiner-lott}
  \dcref{Claim 2}
  \group{Kleiner--Lott Overview}
 \cite{Perelman2} \label{introclaim2}
There is a well-defined surgery procedure on $M$ producing a smooth
post-surgery manifold $M'$ with a Riemannian metric $g'$.
\end{proposition}

The procedure specifies the part removed near $M-\Omega$ and glues caps to
the resulting boundary.  Reversing surgery restores discarded Case A
components and performs connected sums, possibly with finitely many new
$S^1\times S^2$ and $\R P^3$ factors.  The former arises when surgery does not
disconnect a component, and the latter from twisted line-bundle components
in Case B.  The postsurgery metric evolves until the next blow-up time, when
the procedure is repeated.

\begin{proposition}[Ricci flow with surgery proposition]
  \label{prop:kl-overview-ricci-flow-approach-manifold-geometrization-ricci-flow-surgery-proposition}
  \source{kleiner-lott}
  \dcref{Claim 3}
  \group{Kleiner--Lott Overview}
 \cite{Perelman2} \label{introclaim3}
One can arrange the surgery procedure so that the surgery times do not
accumulate.
\end{proposition}

Thus alternating Ricci flow and surgery defines an evolution for all time,
although the manifold may become empty.

\begin{proposition}[Ricci flow with surgery proposition]
  \label{prop:kl-overview-ricci-flow-approach-manifold-geometrization-ricci-flow-surgery-proposition-4}
  \source{kleiner-lott}
  \dcref{Claim 4}
  \group{Kleiner--Lott Overview}
 \cite{Colding-Minicozzi,Colding-Minicozzi2,Perelman3}
\label{introclaim4}
If the original manifold $M$ is simply-connected then any Ricci flow with
surgery on $M$ becomes extinct in finite time.
\end{proposition}

More generally, if the prime decomposition of $M$ has no aspherical factors,
every Ricci flow with surgery becomes extinct in finite time.  A connected
manifold $X$ is aspherical when $\pi_k(X)=0$ for all $k>1$, equivalently when
its universal cover is contractible.  Claims \ref{introclaim1},
\ref{introclaim2}, \ref{introclaim3}, and \ref{introclaim4} imply that a
simply-connected $M$ is a connected sum of factors $S^1\times S^2$ and
$S^3/\Gamma$.  Van Kampen's theorem removes the former factors and shows
$M\cong S^3$.

\subsection{Outline of the proof of the Geometrization Conjecture}

Without simple connectivity, Claim \ref{introclaim4} is unavailable and the
flow with surgery may continue indefinitely.  The proof of geometrization is
independent of Claim \ref{introclaim4} and also implies the Poincar\'e
Conjecture.

For $M=H^3/\Gamma$, with $\Gamma$ freely acting, discrete, cocompact, and
orientation-preserving, let $g_{hyp}$ have sectional curvature $-1$.  If
$g_0=r_0^2g_{hyp}$, then
$g(t)=(r_0^2+4t)g_{hyp}$ and
$\widehat g(t)=t^{-1}g(t)$ converges to the metric of sectional curvature
$-1/4$, independently of $r_0$.

Let $M_t$ be the postsurgery manifold at time $t$.  Components admitting a
metric of nonnegative scalar curvature become extinct or admit a flat metric,
so assume each component has a point of strictly negative scalar curvature.
Set $\widehat g(t)=t^{-1}g(t)$ and define $\rho(x,t)$ by
\[
\inf_{B(x,\rho)}\Rm=-\rho^{-2},
\]
where $\Rm$ is computed using $\widehat g(t)$.  For $w>0$, define
\begin{equation}
M^+(w,t)\: =\:
\{x\in M_t:\vol(B(x,\rho(x,t)))>w\rho(x,t)^3\}.
\end{equation}
It is possible that $M^+(w,t)=M_t$ or $M^+(w,t)=\emptyset$; for every
$w>0$, this part approaches the $w$-thick part of a manifold of constant
sectional curvature $-1/4$ as $t$ increases.

\begin{proposition}[existence for sectional curvature]
  \label{prop:kl-overview-ricci-flow-approach-manifold-geometrization-existence-sectional-curvature}
  \source{kleiner-lott}
  \dcref{Claim 5}
  \group{Kleiner--Lott Overview}
 \cite{Perelman2} \label{introclaim5}
There is a finite collection $\{(H_i,x_i)\}_{i=1}^k$ of complete pointed
finite-volume $3$-manifolds with constant sectional curvature $-1/4$ and, for
large $t$, a decreasing function $\alpha(t)\ra0$ and maps
\begin{equation}
f_t\: :\bigsqcup_{i=1}^kH_i\:\supset\:
\bigsqcup_{i=1}^k B\left(x_i,\frac{1}{\alpha(t)}\right)\longrightarrow M_t
\end{equation}
such that
1. $f_t$ is $\alpha(t)$-close to an isometry;
2. the image of $f_t$ contains $M^+(\alpha(t),t)$;
3. the image under $f_t$ of every cuspidal torus of $\{H_i\}_{i=1}^k$ is
incompressible in $M_t$.
\end{proposition}

The proof of Claim \ref{introclaim5} uses Hamilton's earlier work
\cite{Hamiltonn}.

\begin{proposition}[graph manifolds proposition]
  \label{prop:kl-overview-ricci-flow-approach-manifold-geometrization-graph-manifolds-proposition}
  \source{kleiner-lott}
  \dcref{Claim 6}
  \group{Kleiner--Lott Overview}
 \cite{Perelman2,Shioya-Yamaguchi} \label{introclaim6}
Let $Y_t$ be the truncation of $\bigcup_{i=1}^kH_i$ obtained by removing
horoballs at distance approximately $1/(2\alpha(t))$ from $x_i$.  For large
$t$, $M_t-f_t(Y_t)$ is a graph manifold.
\end{proposition}

Claim \ref{introclaim6} is a Riemannian statement about $3$-manifolds that
are locally volume-collapsed with a lower sectional-curvature bound.  Claims
\ref{introclaim5} and \ref{introclaim6}, together with Claims
\ref{introclaim1}-\ref{introclaim3}, imply geometrization (Appendix
\ref{geom}).

\subsection{Claim \ref{introclaim1} and the structure of singularities}

Near points of large scalar curvature, Ricci flow is modeled at scale
$R(x,t)^{-1/2}$ by a $\kappa$-solution.

\begin{proposition}[existence for kappa solutions]
  \label{prop:kl-overview-ricci-flow-approach-manifold-geometrization-existence-kappa-solutions}
  \source{kleiner-lott}
  \dcref{Claim 7}
  \group{Kleiner--Lott Overview}
\label{introclaim7}
Suppose that we have a given Ricci flow solution on a finite time interval.
If $x\in M$ and the scalar curvature $R(x,t)$ is large, then in the time-$t$
slice there is a ball centered at $x$ of radius comparable to
$R(x,t)^{-1/2}$ whose geometry is close to that of a ball in a
$\kappa$-solution.
\end{proposition}

Here $R(x,t)^{-1/2}$ is the curvature scale.  A $\kappa$-solution is an
ancient Ricci flow, defined on a time interval of the form $(-\infty,a)$,
with nonnegative curvature operator on each time slice,
bounded sectional curvature on compact time intervals, and
$\kappa$-noncollapsing at all scales.

\begin{proposition}[rigidity for kappa solutions]
  \label{prop:kl-overview-ricci-flow-approach-manifold-geometrization-rigidity-kappa-solutions}
  \source{kleiner-lott}
  \dcref{Claim 8}
  \group{Kleiner--Lott Overview}
 \cite{Perelman2} \label{introclaim8}
Any three-dimensional oriented $\kappa$-solution $(M_\infty,g_\infty(\cdot))$
falls into one of the following types:
\begin{enumerate}
\item A finite isometric quotient of the round shrinking $3$-sphere.
\item A Ricci flow on a manifold diffeomorphic to $S^3$ or $\R P^3$.
\item A standard shrinking round neck on $\R\times S^2$.
\item A Ricci flow on a manifold diffeomorphic to $\R^3$, with each time
slice asymptotically necklike at infinity.
\item The $\Z_2$-quotient $\R\times_{\Z_2}S^2$ of a shrinking round neck.
\end{enumerate}
\end{proposition}

Claims \ref{introclaim7} and \ref{introclaim8} give, near large scalar
curvature, closed known-topology regions, necks, $3$-ball caps, or twisted
line-bundle caps.  Their neck overlaps assemble into the components described
in Claim \ref{introclaim1}.

For spacetime points $(x_i,t_i)$ with $R(x_i,t_i)\ra\infty$, rescale about
$(x_i,t_i)$ by defining
\begin{equation}
\overline{g}_i(s)\: =\:R(x_i,t_i)g\left(R(x_i,t_i)^{-1}s+t_i\right).
\end{equation}
The new time interval is $[-R(x_i,t_i)t_i,0]$.  A pointed smooth limit is
ancient because
$\lim_{i\rightarrow\infty}-R(x_i,t_i)t_i=-\infty$, and Hamilton--Ivey
pinching implies that it has nonnegative sectional curvature
(Appendix \ref{phiappendix}).

Hamilton's compactness theorem \cite{Hamiltonnn} applies when, for every
$r>0$ and all sufficiently large $i$, sectional curvature is uniformly
bounded on points $(x,s)$ with
$x\in B_0(x_i,r)$ and $s\in[-r^2,0]$, and the time-zero injectivity radii
$\inj(x_i,0)$ have a uniform positive lower bound.  Under the curvature
condition, the latter is equivalent to the volume bound
\begin{equation} \label{bound}
r^{-3}\vol(B_t(x,r))\geq\kappa\:>\:0,
\end{equation}
for metric $r$-balls satisfying $|\Rm|\leq r^{-2}$, uniformly for
$t\in[0,T)$ and relevant $r<\rho$.  Perelman's no local collapsing theorem
\cite{Perelman} supplies suitable $\kappa$ and $\rho$ for every finite-time
Ricci flow, in arbitrary dimension.  The rescaled limit is therefore
$\kappa$-noncollapsed at all scales; in dimension three it is a
$\kappa$-solution.  The remaining curvature estimates yield Claim
\ref{introclaim7}.

\subsection{The proof of Claims \ref{introclaim2} and \ref{introclaim3}}

Claim \ref{introclaim1} supplies the nonsingular limit on $\Omega$ and the
regions removed in surgery.  Their boundaries have nearly round $2$-spheres
and nearly cylindrical collars, allowing standard $3$-ball caps to be glued
with a partition of unity.  Reapplying Claim \ref{introclaim1} after each
postsurgery flow gives successive surgeries.  To prevent accumulation, the
procedure is chosen so that surgery at time $t$ occurs at a definite scale
$h(t)$ and removes a definite volume $v(t)$.  A quantitative version of Claim
\ref{introclaim1}$,$ uniform over prior surgeries but allowed to depend on
$T'$ and the initial metric, provides this scale.  The required no local
collapsing theorem for Ricci flows with surgery is part of the argument for
Claim \ref{introclaim3}.

\section{Overview of
 \emph{The Entropy Formula for the Ricci
Flow and its Geometric
Applications}
\cite{Perelman}}
\label{overview1}

The paper \cite{Perelman} treats nonsingular Ricci flows; surgery is treated
in \cite{Perelman2}.  Sections I.1--I.10 and the first part of I.11 concern
$n$-dimensional flows, while the second part of I.11 and I.12--I.13 concern
three dimensions.  The main conclusion is geometrization for a compact
$3$-manifold that is the initial manifold of a smooth Ricci flow; Hamilton's
earlier result \cite{Hamiltonn} assumed sectional curvature $O(t^{-1})$ as
$t\rightarrow\infty$.

\subsection{Entropy, breathers, and no local collapsing}
\dcref{I.1--I.6}

Let $M$ be a closed $n$-manifold.  A functional $F(g)$ is monotonic under
Ricci flow when $F(g(t))$ is nondecreasing.  Sections I.1--I.2 introduce
${\mathcal F}(g,f)$: if $g(\cdot)$ is Ricci flow and $e^{-f(\cdot)}$ solves
the conjugate heat equation, then ${\mathcal F}(g(t),f(t))$ is nondecreasing,
with equality exactly for gradient steady solitons.  Minimizing over
$\int_Me^{-f}\,dV=1$ gives the monotonic quantity $\lambda_1(g)$, the lowest
eigenvalue of $-4\triangle+R$.

Section I.3 introduces ${\mathcal W}(g,f,\tau)$, which is nondecreasing when
$\tau=t_0-t$ and $(4\pi\tau)^{-n/2}e^{-f}$ solves the conjugate heat equation;
it is constant exactly for a gradient shrinking soliton terminating at $t_0$.
Section I.4 gives: for a finite-time Ricci flow and any scale $\rho>0$,
there is $\kappa>0$ such that
$|\Rm|\leq r^{-2}$ on $B_t(x,r)$ with $r<\rho$ implies
$\vol(B_t(x,r))\geq\kappa r^n$.  Sections I.5--I.6 develop the entropy
identities and motivate reduced volume.

\subsection{Reduced geometry, Harnack estimates, and pseudolocality}
\dcref{I.7--I.10}

For a spacetime point $(p,t_0)$, set $\tau=t_0-t$.  If
$\gamma(0)=p$ for $0\leq\tau\leq\overline\tau$, define
\begin{equation}
{\mathcal L}(\gamma)=\int_0^{\overline\tau}\sqrt{\tau}\,
\left(|\dot\gamma(\tau)|^2+R(\gamma(\tau),t_0-\tau)\right)\,d\tau.
\end{equation}
Let $L(q,\overline\tau)$ be the infimum of ${\mathcal L}$ over such curves
ending at $q$, and set
$l(q,\overline\tau)=L(q,\overline\tau)/(2\sqrt{\overline\tau})$.  The
reduced volume is
\[
\widetilde V(\overline\tau)=\int_M\overline\tau^{-n/2}
e^{-l(q,\overline\tau)}\,\dvol(q).
\]
It is nonincreasing in $\overline\tau$ and constant exactly for a gradient
shrinking soliton terminating at $t_0$.  For each $\overline\tau$ there is
$q(\overline\tau)$ with $l(q(\overline\tau),\overline\tau)\leq n/2$.

The reduced-volume no local collapsing theorem states that for a flow on
$[0,t_0]$ and $\rho>0$, some $\kappa>0$ satisfies
$\vol(B_{t_0}(p,r))\geq\kappa r^n$ whenever $r<\rho$ and
$|\Rm|\leq r^{-2}$ on
\[
\{(x,t):\dist_{t_0}(x,p)\leq r, t_0-r^2\leq t\leq t_0\}.
\]
Here $\kappa$ depends on $\rho,n,t_0$ and bounds for $g(0)$.  Theorem
I.8.2 extends local $\kappa$-noncollapsing: for every $A<\infty$ there is
$\kappa=\kappa(A)>0$ such that, if the flow exists on $[0,r_0^2]$,
$|\Rm(x,t)|\leq r_0^{-2}$ when $\dist_0(x,x_0)<r_0$, and
$\vol(B_0(x_0,r_0))\geq A^{-1}r_0^n$, then the time-$r_0^2$ metric is not
$\kappa$-collapsed below scale $r_0$ at points with
$\dist_{r_0^2}(x,x_0)\leq Ar_0$.

Section I.9 localizes ${\mathcal W}$.  Section I.10 proves pseudolocality:
an initial ball with a scalar-curvature lower bound and subball isoperimetric
constants close to Euclidean has quantitatively bounded curvature nearby at
later times.

\subsection{$\kappa$-solutions}
\dcref{I.11}

A $\kappa$-solution is a nonflat ancient solution that is
$\kappa$-noncollapsed at all scales and has bounded nonnegative curvature
operator on each time slice.  In dimension three, finite-time blowups and
rescalings of high-scalar-curvature regions produce $\kappa$-solutions.
Hamilton's Harnack inequality and comparison geometry imply that their time
slices have vanishing asymptotic volume ratio.  After normalizing
$R(p,0)=1$, pointed $n$-dimensional $\kappa$-solutions have convergent
subsequences; in dimension three the normalized pointed space is compact.
Compact $3$-dimensional $\kappa$-solutions are spherical space forms.  An
oriented noncompact one is diffeomorphic to $\R^3$, or isometric to the round
shrinking cylinder $\R\times S^2$ or its $\Z_2$-quotient
$\R\times_{\Z_2}S^2$.  For every $\epsilon>0$, each noncompact time slice is,
outside a compact set, after rescaling at $y$ by $R(y,t)$,
$\epsilon$-close to
$[-1/\epsilon,1/\epsilon]\times S^2$ with scalar curvature one.

\subsection{Canonical neighborhoods and thick--thin decomposition}
\dcref{I.12--I.13}

Assume the pinching condition
$\Rm\geq-\Phi(R)$, where
$\lim_{s\rightarrow\infty}\Phi(s)/s=0$, and $\kappa$-noncollapsing below
scale one.  Given $\epsilon,\kappa>0$, there is $r_0>0$ such that if
$t_0\geq1$ and $Q=R(x_0,t_0)\geq r_0^{-2}$, then after scaling by $Q$ the
region
\[
\{(x,t):\dist_{t_0}^2(x,x_0)\leq(\epsilon Q)^{-1},
t_0-(\epsilon Q)^{-1}\leq t\leq t_0\}
\]
is $\epsilon$-close to a corresponding subset of a $\kappa$-solution.  This
is Theorem I.12.1; the pinching condition follows from Hamilton--Ivey
pinching (Appendix \ref{phiappendix}).

Theorem I.12.2 gives forward curvature control: for every $A<\infty$ there
are $K=K(A)<\infty$ and $\rho=\rho(A)<\infty$ such that, under the
hypotheses of Theorem I.8.2 and
$r_0^2\Phi(r_0^{-2})<\rho$,
$R(x,r_0^2)\leq Kr_0^{-2}$ on the time-$r_0^2$ ball of radius $Ar_0$
about $x_0$.  Theorem I.12.3 gives backward control: for $w>0$ there are
$\tau=\tau(w)>0$, $\rho=\rho(w)>0$, and $K=K(w)<\infty$ such that if
$\vol(B(x_0,r_0))\geq wr_0^3$ and $\Rm\geq-r_0^{-2}$ on the time-$t_0$
ball, then
\[
R(x,t)<Kr_0^{-2},\qquad
t\in[t_0-\tau r_0^2,t_0],\quad
\dist_t(x,x_0)<\tfrac14r_0,
\]
provided $\phi(r_0^{-2})<\rho$.  Together with pinching, these estimates
show that balls with infimal sectional curvature $-r^{-2}$ are locally
volume-collapsed relative to the lower curvature bound.

For a nonsingular flow, set $\widehat g(t)=t^{-1}g(t)$.  For large $t$ it
satisfies a universal $\Phi$-pinching condition.  Let $\widehat r(x,t)$ be
defined by
\[
\inf\Rm|_{B_t(x,\widehat r)}=-\widehat r^{-2},
\]
and define
\[
M_{thin}(w,t)=\{x:\vol(B_t(x,\widehat r(x,t)))<w\widehat r(x,t)^3\},
\qquad M_{thick}(w,t)=M-M_{thin}(w,t).
\]
For large $t$, points of $M_{thick}(w,t)$ have constants
$\overline\rho=\overline\rho(w)>0$ and $K=K(w)<\infty$ such that
$|\Rm|\leq Kt^{-1}$ on $B(x,\overline\rho\sqrt t)$ and
\[
\vol(B(x,\overline\rho\sqrt t))
\geq \tfrac1{10}w(\overline\rho\sqrt t)^3.
\]
For sequences $t\rightarrow\infty$ and $w\rightarrow0$, the thick part
converges to a complete finite-volume manifold of constant sectional
curvature $-1/4$ with incompressible cuspidal tori, while the thin part is a
graph manifold.  This gives geometrization for the nonsingular flow; the
surgery version is supplied by Claims \ref{introclaim5} and
\ref{introclaim6}.
